\documentclass[12pt]{article}
\usepackage[utf8]{inputenc}
\usepackage[T1]{fontenc}
\usepackage{amsmath,amssymb}
\usepackage{geometry}
\geometry{margin=2.5cm}

\begin{document}

\begin{center}
\fbox{\textbf{ALGEBRA}}\\[2pt]
\fbox{\textbf{LECTURE 8}}
\end{center}
\begin{flushright}
\fbox{\textit{22-XI-1994}}
\end{flushright}

\bigskip
\noindent\underline{\textbf{The group $S_n$}}

\bigskip
\noindent $A=\{1,2,\ldots,n\}$ \quad $S_n \ni \sigma: A\to A$ \quad $\sigma = \begin{pmatrix}1&2&\cdots&n\\ \sigma(1)&\sigma(2)&\cdots&\sigma(n)\end{pmatrix}$

\noindent $|S_n|=n!$ \qquad $A\xrightarrow{\ \pi\ }A\xrightarrow{\ \sigma\ }A.$

\bigskip
\noindent $(S_n,\circ,e)$ is a group.

\bigskip
\noindent\fbox{Def}: \ Let $p\in\mathbb{N}$, $1<p\leq n$. \ A permutation $\alpha\in S_n$ is called a \underline{cyclic permutation}

\noindent if $\exists\,i_1,\ldots,i_p\in\{1,2,\ldots,n\}$ distinct, with

\[
\alpha(i_1)=i_2, \ \alpha(i_2)=i_3, \ \ldots, \ \alpha(i_{p-1})=i_p, \ \alpha(i_p)=i_1
\]

\bigskip
\noindent A 2-cycle is called a \underline{transposition}.

\bigskip
\noindent\fbox{Prop} \ If $\alpha\in S_n$ is a $p$-cycle, then ord$(\alpha)=p \ \Leftrightarrow \ \begin{cases}\alpha^p=e\\ \alpha^i\neq e \ (\forall)\,i=\overline{1,p-1}\end{cases}$

\bigskip
\noindent In particular, $\tau\in S_n$, $\tau=(i,j) \Rightarrow o(\tau)=2$.

\bigskip
\noindent\underline{Proof}:

\[
\tau\big(\tau(k)\big)=
\begin{cases}
j, & k=i\\
i, & k=j\\
k, & k\neq i,\ k\neq j
\end{cases}
\qquad k\in\{1,2,\ldots,n\}
\]

\bigskip
\noindent $(\alpha^p)(i_1) = \underbrace{\alpha\big(\alpha(\cdots\alpha(i_1))\big)}_{p \text{ times}} = e.$

\bigskip
\noindent\underline{Remark}: \ In $S_n$ there are $n!$ permutations and $C_n^2$ transpositions,

\bigskip
\noindent\fbox{T} \ $(\forall)\,\sigma\in S_n$ \ $(n>1) \Rightarrow \exists\,\tau_1,\ldots,\tau_m\in S_n$ such that $\sigma=\tau_1\circ\cdots\circ\tau_m$.

\bigskip
\noindent\fbox{Proof} \ Given $\sigma\in S_n$, \ $d_\sigma := \text{card}\{i \mid 0<i\leq n,\ \sigma(i)\neq i\}$

\[
0\leq d_\sigma\leq n.
\]

\bigskip
\noindent If $d_\sigma=0 \Rightarrow \sigma=1=e=\tau\circ\tau=\tau\circ\tau\circ\tau\circ\tau=\tau'\circ\tau'\circ\tau''\circ\tau''\circ\tau'''\circ\tau'''$

\bigskip
\noindent Induction on $d_\sigma$.

\noindent Induction hypothesis\footnote{the exact label on this line (``1o.'' in the original) is unclear.} $d_\sigma>0$.

\bigskip
\noindent Claim: it holds for $\pi\in S_n$ with $d_\pi<d_\sigma$.

\noindent We show that $\sigma$ can be written as a finite product of transpositions.

\noindent $d_\sigma>0$ (by hypothesis) $\Rightarrow \exists\,i\in\{1,\ldots,n\}$ such that $\sigma(i)\neq i$

\end{document}
